\chapter{Reprodutibilidade dos procedimentos}\label{ap:repro}

Este apêndice reúne os comandos utilizados na validação, na mesma ordem em
que foram executados. Todas as saídas originais estão preservadas em
\texttt{evidencias/}.

\section{Preparação do servidor MCP}

\begin{verbatim}
cd KiCAD-MCP-Server
npm ci
npm run build
python -m venv .venv
source .venv/bin/activate
pip install -r requirements.txt -r requirements-dev.txt
\end{verbatim}

\section{Ambiente do KiCad (instalação snap, Linux)}

\begin{verbatim}
export PYTHONPATH=/snap/kicad/22/usr/lib/python3/dist-packages
export LD_LIBRARY_PATH=/snap/kicad/22/usr/lib/x86_64-linux-gnu:...:/snap/kicad/22/lib
export PATH=/snap/kicad/22/usr/bin:/snap/kicad/22/bin:$PATH
\end{verbatim}

O \emph{launcher} do projeto (\texttt{run-kicad-mcp.sh}) exporta exatamente
essas variáveis antes de iniciar o servidor Node.

\section{Execução das suítes de teste}

\begin{verbatim}
npm run test:ts                # Vitest (tests-ts/)
python -m pytest tests/ -q -o addopts=""   # pytest; addopts original usa pytest-cov
\end{verbatim}

Os resultados completos foram redirecionados para
\texttt{evidencias/testes-vitest-2026-09-10.log} e
\texttt{evidencias/testes-pytest-2026-09-10.log}.

\section{Verificação de regras e geração de figuras}

\begin{verbatim}
kicad-cli pcb drc --output drc.json --format json --severity-all \
    project-kicad/seguidor_de_linhas_4camada2-2026-1.kicad_pcb
kicad-cli pcb render --output render_top.png --side top \
    --width 2200 --height 1600 --quality high --perspective \
    project-kicad/seguidor_de_linhas_4camada2-2026-1.kicad_pcb
kicad-cli pcb render --output render_bottom.png --side bottom \
    --width 2200 --height 1600 --quality high \
    project-kicad/seguidor_de_linhas_4camada2-2026-1.kicad_pcb
kicad-cli sch export pdf -o sch_root.pdf \
    project-kicad/seguidor_de_linhas_4camada2-2026-1.kicad_sch
pdftoppm -png -r 100 -f 1 -l 1 sch_root.pdf sch_root
\end{verbatim}

\section{Contagens de caracterização}

\begin{verbatim}
grep -c "(footprint " seguidor_de_linhas_4camada2-2026-1.kicad_pcb   # 97
grep -oP '\(net \d+' seguidor_de_linhas_4camada2-2026-1.kicad_pcb \
    | sort -u | wc -l                                              # 83
\end{verbatim}

As dimensoes do contorno foram obtidas pela caixa envolvente dos elementos
da camada Edge.Cuts do arquivo \texttt{.kicad\_pcb}.

\section{Compilação dos documentos}

\begin{verbatim}
cd artigo      && ./build.sh   # gera artigo-mcp-kicad-vscode.pdf
cd monografia  && ./build.sh   # gera monografia-mcp-kicad-vscode.pdf
\end{verbatim}

Ambos os \emph{scripts} executam \texttt{pdflatex} e \texttt{bibtex} em
sequencia (pdflatex, bibtex, pdflatex, pdflatex).
